% Testkorpus: kurze Anleitung mit Bild, Querverweis, Formeln und Listen.
\documentclass[a4paper,11pt]{article}
\usepackage{fontspec}
\usepackage[ngerman]{babel}
\usepackage[margin=2.5cm]{geometry}
\usepackage{graphicx}
\usepackage{booktabs}
\usepackage{amsmath,amssymb}
\usepackage{xcolor}
\usepackage{enumitem}
\usepackage{csquotes}
\usepackage{caption}
\usepackage{float}
\usepackage{microtype}
\usepackage{hyperref}
\begin{document}
\tableofcontents
\section{Montage}\label{sec:montage}
Siehe Abbildung~\ref{fig:bild} und Abschnitt~\ref{sec:montage}. \enquote{Zitat} mit Umlauten äöü ß.
\begin{figure}[H]
  \centering
  \rule{4cm}{2cm}
  \caption{Platzhalter}\label{fig:bild}
\end{figure}
\begin{enumerate}[label=\arabic*.]
  \item Schritt $F = m \cdot a$
  \item \textcolor{blue}{Hinweis} mit $\sum_{i=1}^n i = \frac{n(n+1)}{2}$
\end{enumerate}
\begin{align}
  E &= mc^2 \\
  \mathbb{R} &\ni x
\end{align}
\end{document}
